\chapter{参数化积分}
\label{chap:parametric-integrals}

曲线和曲面通常不是某个固定坐标轴方向上的函数图像。参数化把它们放回
欧氏参数域，但积分值不应依赖参数速度或坐标选择。本章从弧长和面积的局部
线性近似出发，建立两类曲线积分、两类曲面积分及其重新参数化规律。

\section{正则曲线与弧长}

\begin{definition}
映射 $\gamma:[a,b]\to\R^n$ 称为 $C^1$ 正则曲线，如果
$\gamma\in C^1$ 且 $\gamma'(t)\ne0$。允许有限分点处导数不存在而每一段
正则时，称为分段 $C^1$ 正则曲线。
\end{definition}

对分割 $P:a=t_0<\cdots<t_m=b$，折线长度为
$\sum_i\norm{\gamma(t_i)-\gamma(t_{i-1})}$；弧长定义为这些长度的上确界。

\begin{theorem}[弧长公式]
\label{thm:arc-length-parametric}
若 $\gamma\in C^1([a,b],\R^n)$，则曲线可求长且
\[
 L(\gamma)=\int_a^b\norm{\gamma'(t)}\,dt.
\]
\end{theorem}

\begin{proof}
由微积分基本定理和三角不等式，任意分割都有
\[
 \sum_i\norm{\gamma(t_i)-\gamma(t_{i-1})}
 \le\sum_i\int_{t_{i-1}}^{t_i}\norm{\gamma'(t)}dt.
\]
反向估计利用 $\gamma'$ 的一致连续性。在每段取 $\xi_i$，则
\[
 \gamma(t_i)-\gamma(t_{i-1})
 =\gamma'(\xi_i)\Delta t_i+r_i,
 \qquad \sum_i\norm{r_i}\le\omega_{\gamma'}(\norm P)(b-a).
\]
所以折线长度至少为相应 Riemann 和减去趋零误差。取细分割并取上确界，
得到反向不等式。
\end{proof}

\begin{theorem}[曲线的重新参数化]
\label{thm:curve-reparametrization}
设 $\phi:[c,d]\to[a,b]$ 是 $C^1$ 微分同胚，并令
$\widetilde\gamma=\gamma\circ\phi$。则
\[
 L(\widetilde\gamma)=L(\gamma).
\]
对连续标量函数 $f$，第一型曲线积分在此变换下不变；对连续向量场 $F$，
第二型曲线积分在 $\phi$ 递增时不变，在 $\phi$ 递减时变号。
\end{theorem}

\begin{proof}
链式法则给出
$\widetilde\gamma'(s)=\gamma'(\phi(s))\phi'(s)$。因此
\[
 \int_c^d\norm{\widetilde\gamma'(s)}\,ds
 =\int_c^d\norm{\gamma'(\phi(s))}|\,\phi'(s)|\,ds
 =\int_a^b\norm{\gamma'(t)}\,dt.
\]
第一型积分只是在被积式中再乘 $f(\gamma(\phi(s)))$，同一变量代换仍适用。
第二型积分中的 $\phi'(s)$ 不取绝对值，故其符号记录参数方向。严格单调
保证曲线的每一段只经过一次；若允许来回折返，积分会按经过次数重复计算。
\end{proof}

\section{两类曲线积分}

对连续标量函数 $f$，定义第一型曲线积分
\[
 \int_\gamma f\,ds
 :=\int_a^bf(\gamma(t))\norm{\gamma'(t)}\,dt.
\]
它不随定向改变。对连续向量场 $F$，定义第二型曲线积分
\[
 \int_\gamma F\cdot d x
 :=\int_a^b\langle F(\gamma(t)),\gamma'(t)\rangle dt.
\]
它在反向参数化下变号。若写 $F=(P_1,\ldots,P_n)$，也记为
$\int_\gamma\sum_jP_j\,dx_j$。

若 $F=\nabla\varphi$，链式法则给出
\[
 \int_\gamma F\cdot dx
 =\varphi(\gamma(b))-\varphi(\gamma(a)).
\]
这说明梯度场的积分只依赖端点。反向命题需要定义域的连通和闭路积分条件，
将在下一章说明。

\begin{example}
取单位圆的正向参数 $\gamma(t)=(\cos t,\sin t)$，$0\le t\le2\pi$。
对 $F(x,y)=(-y,x)$，
\[
 F(\gamma(t))\cdot\gamma'(t)=1,
 \qquad \int_\gamma F\cdot dx=2\pi.
\]
反向绕行使积分变为 $-2\pi$；而 $\int_\gamma1\,ds=2\pi$ 在两种方向下
相同。这个计算同时显示了第一型积分测量几何长度、第二型积分记录定向的
差别。
\end{example}

\section{正则曲面与面积元}

设 $D\subset\R^2$ 为有界 Jordan 区域，$X$ 在 $\overline D$ 的一个开邻域
上属于 $C^1$。若
$X_u\times X_v\ne0$ 在 $\overline D$ 上处处成立，则称 $X|_D$ 为正则
参数化。为避免同一片曲面被重复覆盖，通常还要求 $X$ 在 $D$ 内单射，并在
闭包上除边界识别外保持一一对应。

导数把参数小矩形的两条边映为 $X_u\Delta u$ 和 $X_v\Delta v$，相应
平行四边形面积为
\[
 \norm{X_u\times X_v}\,\Delta u\,\Delta v.
\]
因此定义面积元
\[
 dS=\norm{X_u\times X_v}\,du\,dv.
\]

若 $Y=X\circ\Phi$ 是由平面 $C^1$ 微分同胚重新参数化，则
\[
 Y_s\times Y_t
 =\det D\Phi\,(X_u\times X_v)\circ\Phi.
\]
结合二维变量代换，曲面面积
\[
 \operatorname{area}(S)=\int_D\norm{X_u\times X_v}\,du\,dv
\]
与正则参数无关。

公式中的行列式来自一个直接的双线性计算。若
$\Phi(s,t)=(u(s,t),v(s,t))$，则
\[
 Y_s=X_u u_s+X_v v_s,\qquad
 Y_t=X_u u_t+X_v v_t,
\]
从而叉积中 $X_u\times X_u$ 与 $X_v\times X_v$ 两项消失，剩下
$(u_sv_t-u_tv_s)X_u\times X_v$。绝对值控制面积，行列式的符号则控制
定向。

连续单位法向场 $n$ 的选择称为曲面的定向。一个正则参数给出
$n=(X_u\times X_v)/\norm{X_u\times X_v}$；交换参数次序或使用负
Jacobian 的重新参数化会反向定向。Möbius 带不存在全局连续法向场，说明
并非每个曲面都可定向。

\section{两类曲面积分}

对连续标量函数 $f$，定义
\[
 \int_Sf\,dS
 =\int_Df(X(u,v))\norm{X_u\times X_v}\,du\,dv.
\]
对连续向量场 $F$ 和已定向曲面，通量积分为
\[
 \int_SF\cdot n\,dS
 =\int_D\langle F(X(u,v)),X_u\times X_v\rangle\,du\,dv.
\]
前者与定向无关，后者反向定向时变号。参数化不变性分别由上节的绝对值
公式和有向叉积公式得到。

若曲面是图像 $X(x,y)=(x,y,g(x,y))$，向上定向时
\[
 X_x\times X_y=(-g_x,-g_y,1),\qquad
 dS=\sqrt{1+g_x^2+g_y^2}\,dx\,dy.
\]
因此两类积分可分别写为
\[
 \int_S f\,dS
 =\iint_D f(x,y,g(x,y))\sqrt{1+g_x^2+g_y^2}\,dx\,dy,
\]
\[
 \int_SF\cdot n\,dS
 =\iint_DF(x,y,g(x,y))\cdot(-g_x,-g_y,1)\,dx\,dy.
\]
第二式使用的是有向面积向量而非单位法向；这正好消去了单位化时出现的
根式。
球面参数在两极退化，但两极为面积零的例外；可以删去小极冠后计算，再由
有界性令其面积趋零。与极坐标相同，这种处理必须明确退化集的贡献，而不能
把非正则参数直接称为处处正则。

\section*{习题}
\addcontentsline{toc}{section}{习题}
\subsection*{A组　曲线}
\begin{enumerate}[label=\textbf{\thechapter.\arabic*.}]
 \exerciseitem{arc-length-proof-details}
 补全弧长公式中余项总和的估计。
 \exerciseitem{arc-length-reparametrization}
 证明弧长在严格单调 $C^1$ 重新参数化下不变，并分别处理递增、递减情形。
 \exerciseitem{helix-length-line-integral}
 计算螺旋线 $(\cos t,\sin t,t)$ 在 $0\le t\le2\pi$ 上的长度及
 $f(x,y,z)=z$ 的第一型曲线积分。
 \exerciseitem{orientation-vector-line-integral}
 证明第二型曲线积分在反向曲线下变号，第一型不变。
 \exerciseitem{gradient-fundamental-line-integral}
 用链式法则证明梯度场的端点公式，并计算 $\nabla(x^2y+z)$ 沿任意指定
 端点曲线的积分。
 \exerciseitem{curve-multiple-traversal}
 比较单位圆走一周、两周和先正向后反向的两类曲线积分。
\end{enumerate}
\subsection*{B组　曲面}
\begin{enumerate}[label=\textbf{\thechapter.\arabic*.},resume]
 \exerciseitem{surface-reparametrization-cross-product}
 推导 $Y_s\times Y_t=\det D\Phi\,(X_u\times X_v)\circ\Phi$。
 \exerciseitem{graph-surface-area}
 推导图形曲面的面积元，并计算抛物面 $z=x^2+y^2$ 在单位圆盘上方部分的面积。
 \exerciseitem{sphere-area-parametric}
 用球坐标参数计算半径 $R$ 的球面面积，严格处理两极退化。
 \exerciseitem{sphere-radial-flux}
 计算向量场 $F(x)=x$ 穿过半径 $R$ 球面的外向通量。
 \exerciseitem{surface-orientation-change}
 证明交换参数顺序使通量积分变号而标量曲面积分不变。
 \exerciseitem{cylinder-side-flux}
 参数化圆柱侧面并计算常向量场和径向场的通量。
\end{enumerate}
\subsection*{C组　参数化条件}
\begin{enumerate}[label=\textbf{\thechapter.\arabic*.},resume]
 \exerciseitem{nonregular-curve-warning}
 给出同一线段的一个导数在中点为零的 $C^1$ 参数，说明几何曲线可求长并不
 意味着该参数正则。
 \exerciseitem{surface-double-cover}
 用球面参数的两个角周期说明非单射参数会按重数重复计算面积。
 \optionalexerciseitem{mobius-nonorientable-normal}
 沿 Möbius 带中心线追踪局部法向，证明连续走一周后法向反向，因而不存在
 全局定向。
\end{enumerate}
